\section*{Bab \ref{ch:g10:earth-energy-budget} --- Neraca Energi Bumi}

\begin{solution}{exo:g10:earth-energy-budget:1}
\omterm{def:g10:earth-energy-budget:insolation}{Insolasi} suatu tempat adalah daya cahaya Matahari yang diterima satu meter
persegi mendatar di puncak atmosfer di atasnya (dirata-ratakan atas sehari
atau setahun). \omterm{def:g10:earth-energy-budget:albedo}{Albedo} sebuah permukaan adalah bagian cahaya Matahari yang
jatuh padanya yang dipantulkannya. Salju segar (0.8 sampai 0.9) memantulkan
jauh lebih banyak daripada samudra terbuka (0.06).
\end{solution}

\begin{solution}{exo:g10:earth-energy-budget:2}
Dipantulkan: $100/340 \approx \qty{29}{\percent}$; diserap udara:
$80/340 \approx \qty{24}{\percent}$; diserap permukaan: $160/340
\approx \qty{47}{\percent}$.
\end{solution}

\begin{solution}{exo:g10:earth-energy-budget:3}
Aspal: $800 \times (1 - 0.04) = \qty{768}{W}$. Salju: $800 \times (1 -
0.85) = \qty{120}{W}$, lebih dari enam kali lebih kecil.
\end{solution}

\begin{solution}{exo:g10:earth-energy-budget:4}
Penguapan: \qty{82}{W/m^2}, dibandingkan dengan \qty{21}{W/m^2} untuk
pemanasan udara. Air yang menguap mengambil \omterm{def:g10:earth-energy-budget:latent-heat}{kalor latennya} dari laut; uapnya
naik bersama udara, dan di tempat uap itu mengembun menjadi titik-titik awan,
sering kali berkilometer lebih tinggi, energi itu dikembalikan kepada udara di
sekitarnya.
\end{solution}

\begin{solution}{exo:g10:earth-energy-budget:5}
Pada bulan Desember belahan bumi utara condong menjauhi Matahari: di kutub
Matahari tetap berada di bawah cakrawala sepanjang hari, jadi tidak ada cahaya
Matahari yang datang. Pada titik balik Desember hal ini berlaku di semua tempat
di sebelah kutub dari $90^\circ - 23.44^\circ =
\qty{66.56}{\degree}$ LU.
\end{solution}

\begin{solution}{exo:g10:earth-energy-budget:6}
$1361 \sin 90^\circ = \qty{1361}{W}$; $1361 \sin 30^\circ =
\qty{680.5}{W}$; $1361 \sin 10^\circ \approx \qty{236}{W}$. Setengah $S_0$
memerlukan $\sin h = 0.5$, yaitu $h = \qty{30}{\degree}$.
\end{solution}

\begin{solution}{exo:g10:earth-energy-budget:7}
Dipantulkan 100, diserap udara 80, diserap permukaan 160:
$100 + 80 + 160 = 340$, seluruh cahaya Matahari tercatat. Di permukaan,
$185 = 25 + 160$: yang tidak dipantulkan diserap.
\end{solution}

\begin{solution}{exo:g10:earth-energy-budget:8}
Piringan: $\pi \times (6.371\times10^6)^2 \approx \qty{1.275e14}{m^2}$; daya:
$1361 \times 1.275\times10^{14} \approx \qty{1.74e17}{W}$. Bola: $4\pi
R^2 \approx \qty{5.10e14}{m^2}$. Rata-rata: $1.74\times10^{17}/5.10\times
10^{14} \approx \qty{340}{W/m^2}$, seperempat $S_0$.
\end{solution}

\begin{solution}{exo:g10:earth-energy-budget:9}
$0.7 \times 5.10\times10^{14} \times 3.16\times10^7 \approx
\qty{1.1e22}{J}$ dalam setahun.
\end{solution}

\begin{solution}{exo:g10:earth-energy-budget:10}
Awan memantulkan $100 - 53 = \qty{47}{W/m^2}$ cahaya Matahari lebih banyak,
yang karenanya tidak diserap. Awan menahan $267 - 239 = \qty{28}{W/m^2}$
inframerah. Secara keseluruhan Bumi kehilangan $47 - 28 = \qty{19}{W/m^2}$:
awan, secara keseluruhan, mendinginkannya.
\end{solution}

\begin{solution}{exo:g10:earth-energy-budget:11}
Dari sekitar 10 Mei sampai akhir Juli, yaitu pada bulan Mei, Juni, dan Juli.
Di khatulistiwa Matahari melintas tepat di atas kepala pada tengah hari pada
kedua ekuinoks, bulan Maret dan September; pada titik balik Matahari melintas
\qty{23.44}{\degree} dari zenit, sehingga khatulistiwa menerima sedikit lebih
sedikit pada bulan Juni dan Desember.
\end{solution}

\begin{solution}{exo:g10:earth-energy-budget:12}
Setiap detik: $82/(2.47\times10^6) \approx \qty{3.32e-5}{kg}$. Setiap hari:
$3.32\times10^{-5} \times 86\,400 \approx \qty{2.9}{kg}$, lapisan setebal
sekitar \qty{2.9}{mm}. Setiap tahun: $3.32\times10^{-5} \times 3.16\times10^7 \approx
\qty{1050}{kg}$, sekitar \qty{1.05}{m} air.
\end{solution}

\begin{solution}{exo:g10:earth-energy-budget:13}
$S_0 \sin\delta = S_0/\pi$ memberikan $\sin\delta = 1/\pi \approx 0.318$, jadi
$\delta \approx \qty{18.6}{\degree}$. Lalu $\sin\lambda = 0.318/\sin 23.44^\circ
\approx 0.318/0.398 \approx 0.800$: $\lambda \approx \qty{53.15}{\degree}$ atau
$180^\circ - 53.15^\circ = \qty{126.85}{\degree}$. Di antara keduanya, $\lambda$
bertambah \qty{73.7}{\degree}, yaitu $73.7/360$ tahun, sekitar
\num{75} hari, dari pertengahan Mei sampai akhir Juli.
\end{solution}

\begin{solution}{exo:g10:earth-energy-budget:14}
Es menyerap $1 - 0.6 = 0.40$ cahaya Matahari, samudra $1 - 0.06 =
0.94$. Daya tambahan: $(0.94 - 0.40) \times 300 \times 10^{12} \approx
\qty{1.6e14}{W}$, lebih besar daripada seluruh fotosintesis planet ini. Energi
tambahan itu menghangatkan air, yang mencairkan lebih banyak es di tepi
bongkah-bongkah es, sehingga makin banyak air gelap yang terbuka: perubahan
itu memperkuat dirinya sendiri.
\end{solution}

\begin{solution}{exo:g10:earth-energy-budget:15}
Terdekat: $1361 \times (149.6/147.1)^2 \approx \qty{1408}{W/m^2}$; terjauh:
$1361 \times (149.6/152.1)^2 \approx \qty{1317}{W/m^2}$; selisihnya sekitar
\qty{7}{\percent}. Pada \qty{60}{\degree} LU \omterm{def:g10:earth-energy-budget:insolation}{insolasi} pada bulan Juni dua puluh
kali lebih besar daripada pada bulan Desember: perubahan beberapa persen tidak
mungkin menyebabkannya. Lagi pula, Bumi paling dekat dengan Matahari pada awal
Januari, selama musim dingin belahan bumi utara, dan kedua belahan bumi, yang
jaraknya ke Matahari sama, mengalami musim yang berlawanan: kemiringan sumbu,
bukan jarak, yang menyebabkan musim.
\end{solution}

\begin{solution}{pb:g10:earth-energy-budget:1}
\textbf{1.} $\pi \times (6.371\times10^6)^2 \approx \qty{1.275e14}{m^2}$.

\textbf{2.} $1361 \times 1.275\times10^{14} \approx \qty{1.735e17}{W}$.

\textbf{3.} $4\pi R^2 \approx \qty{5.101e14}{m^2}$.

\textbf{4.} $1.735\times10^{17}/5.101\times10^{14} \approx \qty{340}{W/m^2}$,
yaitu $S_0/4$.

\textbf{5.} Khatulistiwa: \qty{1361}{W}; \qty{60}{\degree} LU: $1361 \sin
30^\circ = \qty{680.5}{W}$.

\textbf{6.} $100/340 \approx 0.294$.

\textbf{7.} $340 - 100 = \qty{240}{W/m^2}$; atas seluruh Bumi, $240 \times
5.101\times10^{14} \approx \qty{1.224e17}{W}$.

\textbf{8.} $53/340 \approx 0.156$. Awan menaikkan \omterm{def:g10:earth-energy-budget:albedo}{albedo} dengan faktor
$100/53 \approx 1.9$: hampir menggandakannya.

\textbf{9.} $(240/5.67\times10^{-8})^{1/4} \approx \qty{255}{K}$, sekitar
\qty{33}{K} di bawah \qty{288}{K} permukaan.

\textbf{10.} $160 + 342 = \qty{502}{W/m^2}$.

\textbf{11.} $82 + 21 + 398 = \qty{501}{W/m^2}$. Permukaan memperoleh sekitar
\qty{1}{W/m^2}: ketidakseimbangan \qty{0.7}{W/m^2}, setelah dibulatkan.

\textbf{12.} $342/160 \approx 2.1$. Udara mengirim kepada permukaan energi
dalam bentuk inframerah lebih dari dua kali cahaya Matahari yang diserapnya:
itulah efek rumah kaca.

\textbf{13.} $82 \times 3.156\times10^7/(2.47\times10^6) \approx
\qty{1050}{kg}$, lapisan setebal sekitar \qty{1.05}{m}.

\textbf{14.} $0.7 \times 5.101\times10^{14} \approx \qty{3.6e14}{W}$.

\textbf{15.} $2805/72.07 \approx \qty{38.9}{kJ}$ per gram karbon.

\textbf{16.} $104.9\times10^9$ ton sama dengan $1.049\times10^{17}$~g;
$1.049\times10^{17} \times 38.9\times10^3 \approx \qty{4.08e21}{J}$.

\textbf{17.} $4.08\times10^{21}/3.156\times10^7 \approx \qty{1.29e14}{W}$.

\textbf{18.} $1.29\times10^{14}/5.101\times10^{14} \approx
\qty{0.25}{W/m^2}$.

\textbf{19.} $0.7/0.25 \approx 2.8$: ketidakseimbangan Bumi hampir tiga kali
daya yang disimpan fotosintesis.

\textbf{20.} $1.29\times10^{14}/1.224\times10^{17} \approx 0.0011$:
fotosintesis menyimpan \textbf{sekitar \qty{0.1}{\percent}} cahaya Matahari
yang diserap Bumi --- seperseribu.
\end{solution}
